\chapter{L\'aszl\'o Lov\'asz \& Avi Wigderson (2021)}

\section{Biographical Sketches}

\subsection{L\'aszl\'o Lov\'asz}

L\'aszl\'o Lov\'asz was born on 9 March 1948 in Budapest, Hungary. He studied at E\"otv\"os Lor\'and University in Budapest, where he completed his PhD in 1971 under the supervision of Tibor Gallai, with a dissertation on graph theory and combinatorics. He held positions at E\"otv\"os Lor\'and University (1971--1983), Yale University (1983--1999), and Microsoft Research (1999--2013). Lov\'asz has received numerous honors, including the Abel Prize in 2021, the Wolf Prize in Mathematics in 1999, the Knuth Prize in 1999, the Kyoto Prize in 2010, and the Fulkerson Prize in 1982 and 2012.

Lov\'asz was a prodigy in combinatorial mathematics. By age 16, he had already published his first research paper, and by his early twenties he was one of the leading graph theorists in the world. His work is characterized by an extraordinary ability to bring sophisticated mathematical tools---from linear algebra, topology, and optimization---to bear on problems in discrete mathematics.

\subsection{Avi Wigderson}

Avi Wigderson was born on 9 September 1956 in Haifa, Israel. He studied at the Technion and Princeton University, where he completed his PhD in 1983 under the supervision of Richard Lipton. He held positions at the Hebrew University of Jerusalem (1983--1999) and the Institute for Advanced Study in Princeton (1999--present). Wigderson has received the Abel Prize in 2021, the Turing Award in 2023, the Wolf Prize in Mathematics in 2014, and the Knuth Prize in 2009.

Wigderson is one of the founding figures of computational complexity theory. His work on interactive proofs, zero-knowledge protocols, and the role of randomness in computation has fundamentally shaped our understanding of what can be efficiently computed and verified.

\section{High-Level View for a General Audience}

\subsection{What Did Lov\'asz and Wigderson Do?}

Imagine you have a massive social network with billions of users. You want to find communities, recommend friends, and detect fraud. The mathematical language for studying networks is \emph{graph theory}.

Lov\'asz transformed graph theory by bringing in powerful methods from algebra, topology, and optimization. He solved problems that had seemed impossible, like computing the Shannon capacity of a graph---a number that measures how much information can be reliably transmitted through a noisy communication channel. His \emph{semidefinite programming} approach to this problem introduced a new paradigm for combinatorial optimization.

Now imagine you are using a computer program that relies on randomness---say, a cryptographic protocol for secure communication. Is the randomness truly necessary, or could the same result be achieved without it? Wigderson showed that randomness can be eliminated from many algorithms, provided we are willing to assume that certain computational problems are hard. This is the ``hardness versus randomness'' paradigm.

\subsection{Why Does It Matter?}

Their work has transformed:
\begin{itemize}
\item \textbf{Computer Science}: The understanding of computation, randomness, and complexity.
\item \textbf{Algorithms}: Efficient algorithms for network analysis and optimization.
\item \textbf{Cryptography}: Zero-knowledge proofs and secure multiparty computation.
\item \textbf{Network Science}: The structure of large networks.
\item \textbf{P vs. NP}: The central open problem of computer science.
\end{itemize}

\section{The Origin of the Problem}

\subsection{The Shannon Capacity of a Graph}

Claude Shannon (1956) asked: given a noisy communication channel, what is the maximum number of distinct messages that can be sent with zero probability of error, per use of the channel?

For a channel where some symbols can be confused, the zero-error capacity is given by the Shannon capacity $\Theta(G)$ of the confusability graph $G$. Here:
\begin{itemize}
\item Vertices = symbols that can be transmitted.
\item Edges connect confusable symbols.
\item The capacity $\Theta(G)$ is the maximum size of a set of non-confusable $n$-letter words, to the $1/n$ power, as $n \to \infty$.
\end{itemize}

Computing $\Theta(G)$ was known to be extremely difficult. For example, the Shannon capacity of the 5-cycle $C_5$ was an open problem for over 20 years until Lov\'asz computed it as $\sqrt{5}$ using his $\vartheta$ function.

\subsection{The P vs. NP Problem and Randomness}

The P vs. NP problem asks whether problems that can be checked quickly can also be solved quickly. Wigderson's work explores the role of randomness in computation. The central question: can randomness make hard problems easy? And conversely, can we eliminate randomness from algorithms without losing efficiency?

The class BPP (bounded-error probabilistic polynomial time) consists of problems that can be solved by randomized algorithms with high probability. The question of whether $P = BPP$---whether every problem solvable with randomness is also solvable without it---was a central open question.

Wigderson's revolutionary insight was that the answer depends on the existence of hard problems: if there exists a problem that requires exponential-size circuits, then randomization is not needed. This ``hardness vs. randomness'' principle has become a cornerstone of computational complexity theory.

\subsection{Probabilistic Proof Systems}

Building on the work of Goldwasser, Micali, and Rackoff on interactive proofs, Wigderson showed that the class IP (problems with interactive proofs) equals PSPACE (problems solvable with polynomial space). This was a stunning result: any problem solvable with polynomial space can be verified by an interactive protocol.

An interactive proof system for a language $L$ is a protocol between a computationally unbounded prover $P$ and a polynomial-time verifier $V$:
\begin{itemize}
\item Completeness: if $x \in L$, the prover can convince the verifier with high probability.
\item Soundness: if $x \notin L$, no prover can convince the verifier except with small probability.
\end{itemize}

The IP = PSPACE theorem says that the class of languages with interactive proofs is exactly the class of languages solvable with polynomial space. The key to the proof is an arithmetic technique for encoding Boolean formulas as polynomials, which allows the verifier to check properties of the formula by evaluating the polynomial at random points.

\subsection{The PCP Theorem and Its Origins}

The PCP theorem (probabilistically checkable proofs) grew out of the theory of interactive proofs and program checking. The theorem states that every NP problem has a proof that can be verified by reading only a constant number of bits of the proof.

More precisely: NP is exactly the class of languages $L$ for which there exists a polynomial-time verifier $V$ such that:
\begin{itemize}
\item If $x \in L$, there exists a proof $\pi$ such that $V(x, \pi)$ accepts with probability 1 after reading $O(1)$ bits of $\pi$.
\item If $x \notin L$, for every proof $\pi$, $V(x, \pi)$ accepts with probability at most $1/2$ after reading $O(1)$ bits of $\pi$.
\end{itemize}

The PCP theorem is one of the deepest results in theoretical computer science, with profound implications for the hardness of approximation.

\section{Deep Insight into the Problem}

\subsection{The $\vartheta$ Function of a Graph}

Lov\'asz introduced the $\vartheta$ function as a semidefinite programming relaxation of graph invariants. For a graph $G = (V,E)$:
\[
\vartheta(G) = \max_{v \mapsto u_v \in S^{n-1}} \sum_{v \in V} \langle c, u_v \rangle^2,
\]
where $c$ is a unit vector and $u_v$ are unit vectors with $\langle u_v, u_w \rangle = 0$ whenever $vw \in E$.

The $\vartheta$ function satisfies:
\[
\alpha(G) \leq \vartheta(G) \leq \overline{\chi}(G),
\]
where $\alpha(G)$ is the independence number and $\overline{\chi}(G)$ is the clique cover number.

For the 5-cycle $C_5$, Lov\'asz computed $\vartheta(C_5) = \sqrt{5}$. Since $\vartheta$ is multiplicative under the strong graph product, this immediately gave $\Theta(C_5) = \sqrt{5}$.

The $\vartheta$ function is the first example of a semidefinite programming relaxation in combinatorial optimization. It can be computed in polynomial time using semidefinite programming, making it a practical tool for bounding the independence number and the Shannon capacity.

\subsection{The Lov\'asz Local Lemma}

The Lov\'asz Local Lemma (LLL) states: if $\mathcal{E}_1, \ldots, \mathcal{E}_n$ are events, each with probability $p$, and each event depends on at most $d$ other events, then if $ep(d+1) < 1$, there is a positive probability that none of the events occur.

More generally, if each event $\mathcal{E}_i$ has probability at most $p_i$ and depends on at most $d$ other events, and if $\sum_{i=1}^n p_i 2^d < 1/2$ (in the symmetric version), then with positive probability none of the events occur.

The LLL is a powerful tool for proving the existence of combinatorial structures without constructing them explicitly. Typical applications include:
\begin{itemize}
\item Hypergraph coloring: coloring the vertices of a hypergraph so that no edge is monochromatic.
\item Latin transversals: finding a transversal in a Latin square.
\item Ramsey numbers: lower bounds for Ramsey numbers.
\item Algorithmic LLL: the constructive version (Moser--Tardos) gives efficient algorithms for finding structures whose existence is guaranteed by the LLL.
\end{itemize}

\subsection{The PCP Theorem and Inapproximability}

The PCP theorem states that every NP problem has a probabilistically checkable proof where the verifier reads only a constant number of bits. This implies that many optimization problems are NP-hard to approximate.

In particular, the PCP theorem implies:
\begin{itemize}
\item MAX-3SAT is NP-hard to approximate within a factor of $7/8 + \epsilon$ for any $\epsilon > 0$.
\item MAX-CUT is NP-hard to approximate within a factor of $16/17 + \epsilon$.
\item The traveling salesman problem is NP-hard to approximate within any constant factor.
\end{itemize}

The proof of the PCP theorem is a tour de force of computational complexity, combining algebraic error-correcting codes, recursive composition of proof systems, and the theory of probabilistically checkable proofs.

\subsection{Hardness vs. Randomness}

Wigderson's framework: if there is a problem that requires exponential-size circuits (a hardness assumption), then randomness can be derandomized: $P = BPP$. This leads to explicit constructions of pseudo-random generators.

The key insight: if a problem $L \in \text{DTIME}(2^{O(n)})$ does not have polynomial-size circuits, then there exists a pseudo-random generator $G: \{0,1\}^{O(\log n)} \to \{0,1\}^n$ such that no polynomial-time algorithm can distinguish the output of $G$ from truly random bits.

This pseudo-random generator can be used to derandomize any BPP algorithm: instead of using truly random bits, the algorithm uses the output of $G$ seeded with all possible $O(\log n)$ seeds, and takes the majority vote.

\section{Biggest Contributions and New Tools}

\subsection{Lov\'asz's Contributions}

\begin{itemize}
\item \textbf{The $\vartheta$ Function}: A semidefinite programming bound for the Shannon capacity of a graph. Lov\'asz showed that $\Theta(C_5) = \sqrt{5}$ using this tool, solving a problem that had been open for 20 years.
\item \textbf{The Lov\'asz Local Lemma}: If events are not too dependent, they can all be avoided. This is one of the most widely used probabilistic method tools. The LLL has found applications in hypergraph coloring, the theory of random graphs, and the analysis of algorithms.
\item \textbf{Graph Limits (Graphons)}: A complete theory of dense graph limits, providing an analytic framework for large networks. A graphon is a symmetric measurable function $W: [0,1]^2 \to [0,1]$ that captures the limiting structure of a sequence of graphs.
\item \textbf{The LLL Algorithm}: With Lenstra and Lenstra, factoring polynomials over $\mathbb{Q}$ in polynomial time. The LLL algorithm has become a fundamental tool in computational number theory and cryptography.
\item \textbf{Perfect Graphs}: The characterization of perfect graphs (the Strong Perfect Graph Theorem, with Chudnovsky, Robertson, Seymour, and Thomas). A graph is perfect if its chromatic number equals the size of its largest clique for every induced subgraph.
\item \textbf{The Kruskal--Katona Theorem}: Lov\'asz gave an elegant proof using algebraic methods.
\item \textbf{Matching Theory}: With Plummer, Lov\'asz wrote the definitive text on matching theory, covering both the classical results and the algebraic approach to matchings.
\item \textbf{The Shannon Capacity of the Pentagon}: The computation $\Theta(C_5) = \sqrt{5}$ is one of the most celebrated results in graph theory.
\end{itemize}

\subsection{Wigderson's Contributions}

\begin{itemize}
\item \textbf{Zero-Knowledge Proofs}: With Goldreich and Micali, proving that every NP problem has a zero-knowledge proof system. A zero-knowledge proof allows a prover to convince a verifier of a statement without revealing any information beyond its truth.
\item \textbf{The PCP Theorem}: With Arora, Lund, Motwani, Sudan, and Szegedy: every NP proof can be encoded so that verification reads only a constant number of bits. This theorem is the foundation of the theory of hardness of approximation.
\item \textbf{Hardness vs. Randomness}: If $P \neq NP$, then $P = BPP$ (with Impagliazzo). This gives explicit derandomization from hardness assumptions.
\item \textbf{Multiparty Computation}: Goldreich, Micali, Wigderson: secure multiparty computation tolerating any number of malicious parties. This result showed that any function can be computed securely by multiple parties who do not trust each other.
\item \textbf{The Zig-Zag Product}: With Reingold and Vadhan: a new operation on graphs that constructs expanders with optimal parameters. The zig-zag product is a graph-theoretic analog of the semidirect product in group theory.
\item \textbf{IP = PSPACE}: With Shamir: interactive proof systems capture all of polynomial space. This was a stunning result that showed the surprising power of interaction in proof systems.
\item \textbf{The Sliding Scale Conjecture}: A conjecture about the power of interactive proofs with different numbers of rounds.
\end{itemize}

\subsection{Graph Theory and Combinatorics}

Lov\'asz made foundational contributions to matching theory (with Plummer), graph coloring, and the theory of random graphs. His work on the strong perfect graph theorem is a landmark in graph theory.

The strong perfect graph theorem characterizes perfect graphs as exactly those graphs that contain no odd hole or odd antihole (induced subgraphs). The proof, over 150 pages long, combines deep structural graph theory with the theory of decompositions.

Lov\'asz also made important contributions to the theory of random graphs, including the Lov\'asz--Wormald conjecture on random regular graphs.

\subsection{Semidefinite Programming and Optimization}

The $\vartheta$ function introduced semidefinite programming (SDP) to combinatorial optimization. SDP relaxations have since become a standard tool for approximating hard optimization problems, including:
\begin{itemize}
\item MAX-CUT: the Goemans--Williamson algorithm uses an SDP relaxation to achieve a 0.878-approximation.
\item The sparsest cut problem: the Arora--Rao--Vazirani algorithm uses an SDP relaxation.
\item Graph coloring and graph partitioning problems.
\item The unique games conjecture and its relationship to SDP hierarchies (the Lasserre hierarchy).
\end{itemize}

\subsection{Expanders and Pseudorandomness}

The zig-zag product of Reingold, Vadhan, and Wigderson constructs explicit expander graphs with optimal degree expansion. An expander graph is a sparse graph that is highly connected: random walks on the graph mix rapidly.

The zig-zag product takes two graphs $G$ and $H$ and produces a new graph $G \circ H$ with the property that if $G$ and $H$ are expanders, then $G \circ H$ is also an expander. The degree of $G \circ H$ is $\deg(H)^2$, independent of the size of $G$.

This construction has been used in:
\begin{itemize}
\item The proof that SL = L (Reingold's theorem): undirected connectivity can be solved in logarithmic space.
\item The construction of asymptotically good error-correcting codes.
\item The theory of pseudorandomness and derandomization.
\end{itemize}

\subsection{Interactive Proofs and Their Power}

The IP = PSPACE theorem of Shamir (building on work of Lund, Fortnow, Karloff, and Nisan, which itself built on Wigderson's work) showed that interactive proofs capture exactly the class PSPACE.

The proof uses an algebraic technique: the Boolean formula is encoded as a low-degree polynomial over a large finite field, and the verifier checks properties of the polynomial by random sampling. This technique, called ``arithmetization,'' has become a fundamental tool in complexity theory.

Interactive proofs have found practical applications:
\begin{itemize}
\item Delegation of computation: allowing a weak client to verify the computation performed by a powerful server.
\item Verifiable computation in blockchain systems.
\item Proofs of work and proofs of stake.
\end{itemize}

\section{Impact of the Discovery}

\subsection{Impact on Computer Science}

The PCP theorem revolutionized the theory of hardness of approximation, showing that many optimization problems (like MAX-3SAT and MAX-CUT) are NP-hard to approximate beyond certain thresholds. Before the PCP theorem, it was known that some optimization problems were NP-hard to solve exactly, but it was not known that they were hard to approximate.

Zero-knowledge proofs, introduced by Goldwasser, Micali, and Wigderson, are foundational for modern cryptography. They allow one party to prove possession of information without revealing the information itself. Practical implementations include:
\begin{itemize}
\item zk-SNARKs (zero-knowledge succinct non-interactive arguments of knowledge).
\item zk-STARKs (zero-knowledge scalable transparent arguments of knowledge).
\item Bulletproofs and other efficient zero-knowledge protocols.
\end{itemize}

The $\vartheta$ function is the prototype for semidefinite programming approaches to combinatorial optimization. The Goemans--Williamson algorithm for MAX-CUT uses a similar SDP relaxation.

\subsection{Impact on Mathematics}

Graph limits (graphons) provide a new perspective on large networks, connecting graph theory to analysis and probability theory. They have been applied to the study of quasirandom graphs, property testing, and network science.

The LLL algorithm is widely used in number theory (factoring polynomials, solving Diophantine equations) and cryptography (breaking knapsack-based cryptosystems).

The strong perfect graph theorem, with Lov\'asz as a key contributor, is one of the deepest results in graph theory.

\subsection{Impact on Cryptography}

Zero-knowledge proofs are used in blockchain technology for privacy-preserving transactions (e.g., zk-SNARKs). Secure multiparty computation is used for privacy-preserving data analysis and electronic auctions.

The combination of zero-knowledge proofs and blockchain technology has created an entire industry of privacy-preserving cryptocurrencies, including Zcash and Monero.

Secure multiparty computation has been used in practice for:
\begin{itemize}
\item Privacy-preserving auctions in Denmark (sugar beet auctions).
\item Privacy-preserving statistical analysis of genomic data.
\item Secure comparison of financial data across institutions.
\end{itemize}

\section{Current Developments}

\subsection{The Unique Games Conjecture}

The Unique Games Conjecture (UGC) of Khot is a central open problem in approximation algorithms. It has deep connections to the $\vartheta$ function and semidefinite programming. If true, the UGC would imply that the Goemans--Williamson algorithm for MAX-CUT is optimal.

The UGC states that for any $\epsilon > 0$, there exists a constraint satisfaction problem with constraints on two variables that are bijections (unique games) such that it is NP-hard to distinguish between instances that are $(1-\epsilon)$-satisfiable and instances that are $\epsilon$-satisfiable.

Recent work has shown that the UGC is equivalent to the existence of a certain type of integrality gap in the Lasserre SDP hierarchy.

\subsection{Post-Quantum Cryptography}

Building cryptography that resists quantum attacks is a major research direction. Lattice-based cryptography, code-based cryptography, and multivariate cryptography are the leading candidates.

The LLL algorithm is actually a threat to some cryptographic systems: it can break knapsack-based cryptosystems and some lattice-based schemes with poor parameter choices. However, modern lattice-based cryptography (like the Learning With Errors problem) is believed to be secure against both classical and quantum attacks.

\subsection{Interactive Proof Systems and Delegation of Computation}

Recent work on interactive proof systems (including the work of Goldwasser, Kalai, and Rothblum on delegating computation) builds on the foundations laid by Wigderson.

The key idea is that a powerful server can convince a weak client that it has correctly performed a computation, using only logarithmic communication and polylogarithmic verification time. This is called ``verifiable computation'' or ``delegation of computation.''

\subsection{Graph Neural Networks}

The theory of graph limits (graphons) has found applications in machine learning, particularly in the analysis of graph neural networks. Graph neural networks are deep learning architectures designed to process graph-structured data.

The connection between graphons and graph neural networks is:
\begin{itemize}
\item Graphons model the limit of a sequence of graphs, capturing the structural properties of large networks.
\item Graph neural networks learn functions on graphs that are continuous with respect to the graphon topology.
\item The theory of graphon approximation provides guarantees for the generalization of graph neural networks.
\end{itemize}

\subsection{The Sliding Scale Conjecture}

The sliding scale conjecture of Wigderson and others relates the power of interactive proofs to the number of rounds of interaction. It states that there exist languages that require a certain number of rounds of interaction in a proof system, forming a ``sliding scale'' of complexity.

\subsection{Algorithmic LLL and the Constructive Lov\'asz Local Lemma}

The Moser--Tardos algorithmic version of the Lov\'asz Local Lemma provides efficient randomized algorithms for finding structures whose existence is guaranteed by the LLL. The algorithm is remarkably simple: repeatedly resample variables that violate constraints, and the expected running time is polynomial.

The algorithmic LLL has been applied to:
\begin{itemize}
\item Graph coloring and hypergraph coloring.
\item Packet routing in networks.
\item Constraint satisfaction problems.
\item The design of efficient algorithms for combinatorial structures.
\end{itemize}

\subsection{Derandomization and Pseudo-Random Generators}

Following the hardness vs. randomness paradigm, there has been extensive work on constructing explicit pseudo-random generators and derandomizing algorithms. Current research focuses on:

\begin{itemize}
\item Unconditional derandomization of specific algorithms.
\item The relationship between circuit lower bounds and derandomization.
\item Average-case hardness and the existence of one-way functions.
\item The theory of randomness extractors and their applications.
\end{itemize}

\subsection{The LLL Algorithm and Lattice Problems}

The LLL algorithm has been refined and extended since its introduction:
\begin{itemize}
\item The BKZ (Block Korkin--Zolotarev) algorithm improves the approximation factor.
\item The LLL algorithm is used in the proof of the polynomial-time equivalence of certain lattice problems.
\item Lattice reduction algorithms are fundamental to modern cryptography, both for attacking and for building cryptographic systems.
\end{itemize}

\section{Conclusion}

Lov\'asz and Wigderson bridged mathematics and computer science, solving deep problems while creating tools essential for modern computing, cryptography, and algorithm design. Lov\'asz introduced geometric and algebraic methods that transformed combinatorics, while Wigderson revealed the deep interplay between randomness, hardness, and proof.

Their work is a testament to the power of interdisciplinary thinking. The methods they developed---semidefinite programming, the probabilistic method, interactive proofs, and pseudorandomness---have become foundational tools across mathematics and computer science.

The problems they opened---the power of approximation, the nature of randomness, and the structure of large networks---continue to drive research at the frontier of computation and mathematics.


\section{Behind the Breakthrough: The Human Story}
The LLL algorithm has been used to break almost every knapsack-based cryptosystem, showing how a pure math lattice algorithm had a direct, devastating impact on computer security.

\section{Pedagogical Metaphors and Lineage}
\subsection{Signature Metaphor}
``Finding short vectors in high-dimensional lattices, and verifying proofs using randomized checks.''

\subsection{Mathematical Lineage}
Lov\'asz advised by Tibor Gallai. Wigderson advised by Richard Lipton.

\section{Applied Science and Computational Algorithms}
\subsection{Key Takeaways for Applied Science}
Lattice reduction is critical to breaking cryptosystems and designing secure post-quantum schemes (Lattice-based cryptography like LWE), and solving mixed-integer linear programming.

\subsection{Algorithms and Numerical Implementations}
The LLL (Lenstra-Lenstra-Lov\'asz) algorithm finds short vectors in lattices in polynomial time, and semidefinite programming solvers compute the Lov\'asz Theta function for graph optimization.

\section{The Research Frontier}
The P vs. NP problem, refining lower bounds in circuit complexity, and arithmetic complexity approaches to graph isomorphism.

\section{Guided Exercises and Challenges}
\begin{enumerate}[label=\textbf{\arabic*.}]
  \item State the Lov\'asz Local Lemma and explain how it allows one to prove the existence of objects satisfying certain constraints.
  \item Define the PCP (Probabilistically Checkable Proofs) theorem and explain its significance for the inapproximability of NP-hard problems.
  \item Explain Wigderson's de-randomization results: under what conditions can randomized algorithms be converted to deterministic ones ($P = BPP$)?
\end{enumerate}

\section{Curriculum Map and Further Reading}
For students wishing to delve deeper into the mathematical machinery underlying these breakthroughs, the following learning path is recommended:
\begin{itemize}
  \item L. Lov\'asz, \emph{Large Networks and Graph Limits}, American Mathematical Society.
  \item Sanjeev Arora and Boaz Barak, \emph{Computational Complexity: A Modern Approach}, Cambridge.
\end{itemize}

\section*{Bibliography}
\begin{enumerate}
\item L. Lov\'asz, ``On the Shannon capacity of a graph,'' IEEE Transactions on Information Theory 25 (1979), 1--7.
\item L. Lov\'asz, ``Large Networks and Graph Limits,'' American Mathematical Society, 2012.
\item L. Lov\'asz, ``The Lov\'asz Local Lemma and its applications,'' in ``Combinatorics, Paul Erd\H{o}s is Eighty,'' Bolyai Society, 1993.
\item A. Wigderson, ``Mathematics and Computation,'' Princeton University Press, 2019.
\item S. Goldwasser, S. Micali, and A. Wigderson, ``How to play any mental game,'' in ``Proceedings of STOC 1987,'' 218--229.
\item S. Arora and B. Barak, ``Computational Complexity: A Modern Approach,'' Cambridge University Press, 2009.
\item S. Arora, C. Lund, R. Motwani, M. Sudan, and M. Szegedy, ``Proof verification and the hardness of approximation problems,'' Journal of the ACM 45 (1998), 501--555.
\item O. Goldreich, ``Foundations of Cryptography,'' Cambridge University Press, 2001.
\item R. Impagliazzo and A. Wigderson, ``$P = BPP$ if $E$ requires exponential circuits,'' in ``Proceedings of STOC 1997,'' 220--229.
\item O. Reingold, S. Vadhan, and A. Wigderson, ``Entropy waves, the zig-zag graph product, and new constant-degree expanders,'' Annals of Mathematics 155 (2002), 157--187.
\item L. Lov\'asz and M. D. Plummer, ``Matching Theory,'' North-Holland, 1986.
\item L. Lov\'asz, ``A semidefinite programming relaxation of the Shannon capacity,'' in ``Proceedings of the ICM 1994.''
\item A. K. Lenstra, H. W. Lenstra, and L. Lov\'asz, ``Factoring polynomials with rational coefficients,'' Mathematische Annalen 261 (1982), 515--534.
\item M. Chudnovsky, N. Robertson, P. Seymour, and R. Thomas, ``The strong perfect graph theorem,'' Annals of Mathematics 164 (2006), 51--229.
\item S. Jukna, ``Extremal Combinatorics,'' Springer, 2011.
\item L. Lov\'asz, ``Semidefinite programs and combinatorial optimization,'' in ``Proceedings of the ICM 1998.''
\item A. Shamir, ``IP = PSPACE,'' Journal of the ACM 39 (1992), 869--877.
\item S. Goldwasser and S. Micali, ``Probabilistic encryption and how to play mental poker keeping secret all partial information,'' in ``Proceedings of STOC 1982,'' 365--377.
\item O. Goldreich, S. Micali, and A. Wigderson, ``Proofs that yield nothing but their validity, or all languages in NP have zero-knowledge proof systems,'' Journal of the ACM 38 (1991), 690--728.
\item S. Arora and S. Safra, ``Probabilistic checking of proofs: a new characterization of NP,'' Journal of the ACM 45 (1998), 70--122.
\item U. Feige, S. Goldwasser, L. Lov\'asz, S. Safra, and M. Szegedy, ``Interactive proofs and the hardness of approximating cliques,'' Journal of the ACM 43 (1996), 268--292.
\item J. Koml\'os and E. Szemer\'edi, ``Limit distribution for the existence of Hamilton circuits in a random graph,'' Discrete Mathematics 43 (1983), 55--63.
\item M. X. Goemans and D. P. Williamson, ``Improved approximation algorithms for maximum cut and satisfiability problems using semidefinite programming,'' Journal of the ACM 42 (1995), 1115--1145.
\item S. Khot, ``On the power of unique 2-prover 1-round games,'' in ``Proceedings of STOC 2002,'' 767--775.
\item S. Khot, G. Kindler, E. Mossel, and R. O'Donnell, ``Optimal inapproximability results for MAX-CUT and other 2-variable CSPs?'' SIAM Journal on Computing 37 (2007), 319--357.
\item R. A. Moser and G. Tardos, ``A constructive proof of the general Lov\'asz Local Lemma,'' Journal of the ACM 57 (2010), 1--15.
\item O. Reingold, ``Undirected connectivity in log-space,'' Journal of the ACM 55 (2008), 1--24.
\item D. A. Spielman and S.-H. Teng, ``Spectral sparsification of graphs,'' SIAM Journal on Computing 40 (2011), 981--1025.
\item L. Lov\'asz and V. T. S\'os, ``Generalized quasirandom graphs,'' Journal of Combinatorial Theory B 98 (2008), 1184--1203.
\item L. Lov\'asz and B. Szegedy, ``Limits of dense graph sequences,'' Journal of Combinatorial Theory B 96 (2006), 933--957.
\end{enumerate}
